\subsection{Principal Findings}

For RQ1, the strongest diagnostic result is the learnability of the four annotated fault domains once anomalous windows are supplied. Stage 2 Macro-F1 was 0.8980 under the fixed chronological protocol and 0.9041 under leave-log-out. The complete Cascade is more demanding because Stage 1 can miss anomalies or introduce false alarms; its results and the Direct Five-Class comparison are therefore reported separately. The lower detector AUPRC under strict isolation limits the transfer of fixed-split performance to new flights.

For RQ2, hierarchical evidence preserves the provenance of each statistical contribution and supports a quantitative perturbation test. Validation-ranked SHAP masking caused a greater diagnostic performance decrease than equal-size random masking under both replacement strategies. Conservation verifies the aggregation implementation, while the weaker External and Global Position semantic agreement shows that model relevance and physical fault-domain agreement are different properties.

For RQ3, the commit-specific graph converts topic evidence into source-tree inspection candidates, and the controlled paired-replay extension provides a bounded test against known mutated modules. Common-support residual processing reduced normal alarms from 11/12 to 3/12 while retaining 10/12 detections on the examined transfer set. This positive result coexists with EKF2 false alarms, delayed or missed INAV detections, and static publisher ties. It supports further investigation of diagnostic evidence under matched replay conditions, without establishing reliable end-to-end software fault localization.
